\documentclass[11pt]{article}
\usepackage[margin=1in]{geometry}
\usepackage{amsmath,amssymb}
\usepackage[T1]{fontenc}
\usepackage{lmodern}
\usepackage{microtype}
\usepackage{xurl}
\usepackage{booktabs}
\usepackage{array}
\usepackage{hyperref}
\setlength{\emergencystretch}{3em}

\title{Empirical Strategy Blueprint\\Fertility and Housing Supply}
\author{Fertility and Housing Supply Project}
\date{February 20, 2026}

\begin{document}
\maketitle

\section*{Purpose}
This note ranks candidate empirical designs and provides an implementation blueprint for the top two. The ranking criterion is practical: high identification credibility with feasible data assembly in finite time.

\section{Ranking summary}
\begin{center}
\begin{tabular}{p{0.58\textwidth}c}
\toprule
Design & Rank \\
\midrule
State middle-housing legalizations and fertility outcomes & 1 \\
NIMBY political turnover (close elections) and fertility via supply & 2 \\
Family-sized unit composition shocks (shift-share IV) & 3 \\
Short-term rental crackdown shocks and fertility & 4 \\
Single-city mega-reform with legal interruptions & 5 \\
\bottomrule
\end{tabular}
\end{center}

\section{Top design 1: Middle-housing legalizations}
\subsection*{Research question}
Do reforms that legalize middle housing (duplex/triplex/fourplex forms) increase fertility among prime childbearing populations by improving access to family-suitable housing?

\subsection*{Estimand}
Dynamic average treatment effect on treated locations (ATT) for age-specific birth rates.

\subsection*{Identification strategy}
Staggered event-study difference-in-differences with treatment intensity.

Baseline event-study:
\[
BirthRate_{it} = \sum_{k\neq -1}\beta_k 1\{EventTime_{it}=k\}+\alpha_i+\delta_t+X_{it}'\Gamma+\varepsilon_{it}.
\]

Intensity extension:
\[
BirthRate_{it}=\beta(Post_{it}\times Exposure_i)+\alpha_i+\delta_t+X_{it}'\Gamma+\varepsilon_{it}.
\]

\subsection*{Data build}
\begin{itemize}
\item Natality outcomes by mother's age from public vital-statistics sources (county-year minimum).
\item Permits and housing production composition (single-family vs multifamily; where possible, unit-size proxies).
\item Price/rent controls and labor-market controls.
\item Policy adoption and effective-implementation timing panel.
\end{itemize}

\subsection*{Main identifying assumptions}
\begin{itemize}
\item In absence of reform, treated and untreated areas follow parallel fertility trends conditional on fixed effects and controls.
\item No confounding policy shocks perfectly collinear with reform timing after controls.
\item Exposure intensity captures meaningful heterogeneity in reform bite.
\end{itemize}

\subsection*{Falsification and robustness}
\begin{itemize}
\item Pre-trend tests and placebo policy dates.
\item Alternative treatment windows (announcement vs implementation).
\item Alternative exposure definitions.
\item Re-estimation excluding major outliers and migration-heavy areas.
\end{itemize}

\subsection*{Strengths and risks}
\textbf{Strengths:} policy relevance, scalable geography, direct link to housing supply.

\textbf{Risks:} first-stage quantity response may be slow; fertility response may lag reform substantially.

\section{Top design 2: NIMBY political turnover design}
\subsection*{Research question}
Does narrowly electing pro-housing local governments shift housing supply and affordability, and through this channel alter fertility outcomes?

\subsection*{Estimand}
Two-step causal object:
\begin{enumerate}
\item effect of pro-housing victory on supply/price outcomes;
\item effect of induced supply changes on fertility.
\end{enumerate}

\subsection*{Identification strategy}
Close-election RD combined with IV-style mediation.

First-stage equation:
\[
Supply_{it}=\pi\,ProHousingWin_{it}+f(Margin_{it})+\eta_i+\tau_t+u_{it}.
\]
Second-stage equation:
\[
BirthRate_{it}=\beta\,\widehat{Supply}_{it}+f(Margin_{it})+\eta_i+\tau_t+e_{it}.
\]

\subsection*{Data build}
\begin{itemize}
\item Local election returns with close-race margins.
\item Candidate/council platform coding for housing stance.
\item Permit and unit-composition outcomes.
\item Price/rent outcomes and fertility outcomes mapped to compatible geography.
\end{itemize}

\subsection*{Main identifying assumptions}
\begin{itemize}
\item Near the cutoff, winners and losers are quasi-random.
\item Platform coding is sufficiently accurate to separate pro-housing from anti-growth outcomes.
\item Exclusion for second stage: election-induced fertility impacts operate primarily through supply/price channels (to be argued carefully).
\end{itemize}

\subsection*{Falsification and robustness}
\begin{itemize}
\item McCrary-style density and covariate-balance checks at the cutoff.
\item Placebo cutoffs and placebo outcomes.
\item Alternative bandwidths and polynomial orders.
\item Robustness to coding uncertainty (drop ambiguous races).
\end{itemize}

\subsection*{Strengths and risks}
\textbf{Strengths:} direct NIMBY political-economy identification; strong conceptual fit.

\textbf{Risks:} heavy data engineering and potential power constraints from close-race sample size.

\section{Why these two are the best portfolio}
These top two designs are complementary:
\begin{itemize}
\item Design 1 delivers broad external validity and cleaner implementation.
\item Design 2 delivers a sharper political-economy test tied directly to your paper's NIMBY channel.
\end{itemize}

Running both in parallel gives a ``safe'' main design and a ``high-upside'' political design.

\section{Operational sequencing}
\subsection*{Phase 1 (fast start)}
\begin{enumerate}
\item Build policy timing panel and core outcome panel.
\item Estimate first-pass event-study for design 1.
\item Verify first stage on permits and composition.
\end{enumerate}

\subsection*{Phase 2 (NIMBY track)}
\begin{enumerate}
\item Build election and coding dataset.
\item Run close-race diagnostics and first stage.
\item Add second-stage fertility mapping where geography allows.
\end{enumerate}

\section*{Bottom line}
For a credible and publishable empirical chapter, begin with middle-housing reforms as the baseline design. In parallel, develop the close-election NIMBY design as the distinctive contribution that links your theory directly to political supply restrictions and demographic outcomes.

\end{document}
